\section{绪论}

\subsection{研究背景}

计算机辅助设计是机械设计、模具设计和产品结构设计中的基础工具。传统 CAD 软件在工业生产中具有较高成熟度，但其操作模式主要围绕几何草图、拉伸旋转、布尔运算、约束求解和参数编辑展开。这种工作方式对专业人员十分高效，但对于初学者和课程设计学生而言学习成本较高。尤其在毕业设计或课堂演示场景中，设计需求往往首先以自然语言、尺寸描述和工艺要求等文本形式出现，随后才被人工翻译为 CAD 操作步骤。

近几年，大语言模型在自然语言理解、信息抽取、结构化输出和代码生成方面取得明显进展。Transformer 架构的提出奠定了现代大模型的发展基础 \cite{vaswani2017attention}，后续通用大模型在复杂语义理解和结构化信息生成上表现出较强能力 \cite{openai2023gpt4}。这为“自然语言到工程表达”的系统设计提供了新的可能。

但在工程场景中，单纯依赖大模型直接生成 CAD 脚本存在两个突出问题：一是脚本语法与几何语义约束强，输出稍有偏差就可能无法运行；二是机械结构需求通常包含尺寸依赖、装配关系和坐标语义，若缺少明确的中间表示，后续生成阶段难以稳定控制。因此，如何在“模型灵活性”和“工程可验证性”之间取得平衡，是本课题需要解决的核心问题。

\subsection{研究意义}

本文研究的意义主要体现在三个方面：其一，能够为机械类课程设计和系统演示提供一条从文本需求到参数化脚本的完整教学链路；其二，通过“自然语言分析 + 强类型 Plan + 确定性模板”的组合式架构，为受限任务域中的 AI 工程软件设计提供了可复用的方法；其三，该系统适合快速原型和概念展示场景，具有明确的工具化落地价值。

\subsection{国内外研究现状}

围绕 CAD 自动建模与结构生成，已有研究在深度学习和程序化建模方向展开探索。SketchGraphs 提供了大规模 CAD 草图与操作数据集，为学习型 CAD 建模研究提供了基础 \cite{seff2020sketchgraphs}。DeepCAD 进一步尝试利用神经网络学习 CAD 序列表示，实现面向 B-Rep 构造过程的生成 \cite{wu2021deepcad}。近年的 Text2CAD 研究则开始关注由文本描述驱动 CAD 表达生成的问题 \cite{khan2024text2cad}。

然而，面向本科毕业设计与课程演示的系统实现，与上述研究型工作存在明显差异。研究型工作通常强调大规模数据集、复杂模型训练和几何重建质量；而教学与原型系统更关心输入是否容易理解、输出是否稳定、错误是否能被尽早发现，以及生成产物是否便于继续人工修改。基于这一差异，本文并不追求复杂几何体或开放域建模，而是将任务域限定为简单装配结构，并采用更强调工程控制的实现路线。

另一方面，Onshape 提供的 FeatureScript 机制为参数化特征开发提供了良好支撑 \cite{onshape_fsdoc,onshape_featurescript}. FeatureScript 本质上是一种面向 CAD 特征定义的脚本化语言，能够通过明确的几何 API 与布尔运算构造三维实体。对于本课题而言，FeatureScript 既具备参数化建模表达能力，又易于通过代码方式生成，因此成为系统的目标输出形式。

\subsection{本文主要工作}

围绕“自然语言驱动简单机械结构建模”这一目标，本文主要完成了以下工作：
\begin{enumerate}[leftmargin=2em]
  \item 设计了一个面向简单机械装配的受限中间表示 \texttt{AssemblyPlan}，用于承载零件形状、尺寸参数、空间位置和装配关系。
  \item 实现了基于大语言模型的需求分析模块，将自然语言机械需求转换为结构化 Plan，并结合强类型校验保证输出稳定性。
  \item 实现了确定性 FeatureScript 模板生成模块，避免让大模型直接自由生成 CAD 脚本，从而提高脚本的可控性和可验证性。
  \item 构建了包含 CLI 与本地 Web UI 的双交互入口，支持分析、生成、全流程构建等多种使用方式。
  \item 基于自动化测试和典型样例完成系统验证，并分析了系统当前能力边界与后续扩展方向。
\end{enumerate}

\subsection{论文结构安排}

本文后续内容安排如下：第 2 章介绍相关技术与系统需求；第 3 章给出系统总体设计；第 4 章说明关键模块实现；第 5 章展示系统测试与结果分析；第 6 章总结全文并展望后续工作。
